\section{Numerical and Temporal Verification}
\label{sec:verification}
\subsection{Separate Oracles for Angles and Pixels}
The ten-suite regression was rerun on October 23, 2026 at source revision
\code{5126c37}. Full hashes, commands, inputs, and logs are archived with the
paper. Results below refer to this run unless a historical scope is stated.

The orientation reference is extracted from the compiled C++ translation
unit, with source and table hashes recorded. A portable C++ harness is
compiled using MSVC with precise floating-point semantics for table
construction; an independently expressed Python model checks the integer
operations. The 79,397 requests compare 245,237 integer outputs, including
inverse-trigonometric boundaries, division, integer square root, signed
extremes, directed camera cases, and 2,000 seeded random frames. All recorded
outputs agree exactly. This establishes agreement with the chosen discrete
reference, not agreement with ideal real-valued rotations or measured rig
orientation.

The pixel oracle is separate. It uses the preserved EO/IR segment maps and
sampling arithmetic, rather than the newer C application's optical profile.
Two renderer suites exercise legacy static input and the frozen gravity-pose
interface. Each compares 14,112,000 pixels exactly across EO and IR cases,
including zero, mixed, and limit-angle configurations. Consequently angle
agreement and pixel agreement are not conflated into an unsupported claim
of end-to-end equality to a different optical pipeline.

\begin{table}[t]
\centering\footnotesize
\caption{Recorded verification scopes and outcomes. Counts refer to tested
cases, not exhaustive coverage of the complete system state space.}
\label{tab:verification}
\begin{tabularx}{\columnwidth}{@{}lY@{}}
\toprule
Scope & Check and observed result \\
\midrule
C/Python arithmetic & 79,397 requests; 245,237 integer outputs agree \\
Camera projection RTL & 1,342 compiled-C vectors; exact outputs \\
Six-camera controller & 266 cases; atomic publication and limiting \\
Pose association & 21 lease cases; independent EO/IR histories, cancellation, fallback \\
Trigger epoch & Phase/reset tests; all 65,536 counter increments \\
I\textsuperscript{2}C interface & 109 wire-level checks; partial/malformed transaction rejection \\
CDC mailbox & Three clock ratios, clock stoppage, reset and coherent transfer \\
Renderer & Two suites, each 14,112,000 matching pixels; eight table cases \\
Capture/selection & 116 envelope checks; 19 single-selection checks; marker-order test \\
\bottomrule
\end{tabularx}
\end{table}

\subsection{State Invariants and Failure Cases}
Table~\ref{tab:verification} complements numeric tests with transaction and
state-machine tests. Publication occurs only after a complete pose burst;
changing input registers while computation is busy cannot change its result.
A source lease selects the corresponding family's history entry, not the
newest pose and not an equal numeric epoch from the other sensor family.
Missing, stale, disabled, overwritten, and calibration-mismatched entries
exercise the static fallback. Mode and lease changes cancel pending work,
while reset clears valid state. The renderer tests vary live inputs during
table generation and replay to check that accepted coefficients remain
constant.

The capture-envelope check has a different character: affine coordinates
within each segment and pitch monotonicity reduce a bounded geometric
coverage problem to endpoint checks over the supported discrete roll range.
This proves coverage of the defined fetch envelope under those model
assumptions. It does not prove that camera capture, DDR arbitration, source
bank ownership, and display publication meet every real-time deadline.
Likewise functional CDC tests do not prove metastability margins; those
require implementation constraints and physical review.

The orientation engine takes 531--601 clock cycles in the camera tests; the
enabled six-camera controller takes 3,556--7,319 cycles, depending on whether
the global rescale needs a second pass. At the nominal 233.4 MHz UI clock,
that controller range corresponds to approximately 15.2--31.4 microseconds.
This is a conversion of simulation cycles, not a measured camera-to-display
latency, and excludes table building, source replay, memory contention,
publication, and scanout. The firmware control and calibration-storage tests
are archived separately; they do not validate the physical filter response.
